% -*- TeX:UK -*-
\documentclass[10pt]{beamer}
\usetheme{metropolis}
%\useinnertheme{rectangles}
\setbeamercovered{%
still covered={\opaqueness<1->{15}},
again covered={\opaqueness<1->{40}}}

\hypersetup{colorlinks,linkcolor=black,urlcolor=brown,citecolor=brown}

\usepackage{amsmath,amssymb,amsthm}
\usepackage{unicode-math}

% We set the Lucida OTF fonts as default
\usepackage{fontspec}
\setmainfont{Lucida Bright OT}
\setsansfont{Lucida Sans OT}
\setmonofont{Lucida Console DK}[Scale=MatchLowercase]

\newfontfamily\webglyphsfont{WebHostingHub-Glyphs}[Scale=0.7]
\newcommand\webglyphs[1]{{\webglyphsfont\symbol{#1}}}
\newcommand\Discussion{\colorbox{white}{\textcolor{black}{\webglyphs{"F134}}}\xspace}
\newcommand\DiscussionI{\colorbox{black}{\textcolor{white}{\webglyphs{"F134}}}\xspace}
\newcommand\DExamples{\colorbox{black}{\textcolor{white}{\webglyphs{"F134} examples?}}}
\newcommand\Reading{\colorbox{black}{\textcolor{white}{\webglyphs{"F0C1}}}\xspace}
\newcommand\ReadingI{\colorbox{white}{\textcolor{black}{\webglyphs{"F0C1}}}\xspace}
\newcommand\Video{\colorbox{white}{\textcolor{black}{\webglyphs{"F03D}}}\xspace}
\newcommand\Attention{\colorbox{black}{\textcolor{orange}{\webglyphs{"F05A}}}\xspace}
\newcommand\HomeWork{\colorbox{white}{\textcolor{black}{\webglyphs{"F5ED}}}\xspace}
\newcommand\HomeWorkI{\colorbox{black}{\textcolor{white}{\webglyphs{"F5ED}}}\xspace}
\newcommand\Advanced{\colorbox{black}{\textcolor{white}{\webglyphs{"F235}}}\xspace}

\newfontfamily\lineabasicfont{linea-basic-10}
\newcommand\basicicons[1]{{\lineabasicfont\symbol{#1}}}
\newcommand\timeforwards{\basicicons{"0079}}
\newcommand\timebackwards{\basicicons{"0064}}

\newfontfamily\lineaweatherfont{linea-weather-10}
\newcommand\weathericons[1]{{\lineaweatherfont\symbol{#1}}}
\newcommand\meteosun{\weathericons{"E038}}
\newcommand\meteosuncloud{\weathericons{"E042}}
\newcommand\meteorain{\weathericons{"E033}}
\newcommand\meteowind{\weathericons{"E054}}

\newfontfamily\uleaffont{Mini Pics Uprooted Leaf}
\newcommand\uleafmpics[1]{{\uleaffont\symbol{#1}}}
\newcommand\lowplants{\uleafmpics{"00CE}}
\newcommand\mediumplant{\uleafmpics{"006A}}
\newcommand\bush{\uleafmpics{"0039}}
\newcommand\smallplant{\uleafmpics{"0030}}
\newcommand\seedling{\uleafmpics{"002F}}
\newcommand\floweringplant{\uleafmpics{"00CA}}

\newfontfamily\utwigfont{Mini Pics Uprooted Twig}
\newcommand\utwigmpics[1]{{\utwigfont\symbol{#1}}}
\newcommand\grassplant{\utwigmpics{"0033}}

\newfontfamily\uinsectfont{Insect Icons}
\newcommand\uinsect[1]{{\uinsectfont\symbol{#1}}}
\newcommand\bug{\uinsect{"006F}}

\usepackage{polyglossia}
\setdefaultlanguage[variant = british, ordinalmonthday = false]{english}

\usepackage[style=authoryear-comp,firstinits,sortcites,maxcitenames=2,%
    mincitenames=1,maxbibnames=10,minbibnames=10,uniquename=mininit,%
    uniquelist=minyear,sortfirstinits=true]{biblatex}
\addbibresource{../references/ecophys.bib}
\renewcommand{\bibfont}{\small}

\usepackage{abbrev}



\begin{document}

\title{PBIO-141\\Sensory and Physiological Ecology\\ of  Plants}
\subtitle{3: Signals, cues, information and evolution}
\author{Pedro J. Aphalo}
\date{January--February 2022}
\institute[Univ.\ of Helsinki]{M.Sc.\ in Plant Biology, University of Helsinki\\[2ex] \url{http://blogs.helsinki.fi/aphalo/}}

  \begin{frame}
    \maketitle
  \end{frame}

  \begin{frame}[c]
    \begin{center}
      \begin{small}
        \copyright 2006--2022 by Pedro J. Aphalo\\
        University of Helsinki, Finland.\\
        \textcolor{blue}{\url{http://blogs.helsinki.fi/senpep-blog/}}\\[2ex]
      \end{small}

      \begin{footnotesize}
        Sensory and Physiological Ecology of Plants slides by Pedro J. Aphalo are licensed under a Creative Commons Attribution-ShareAlike 4.0 International License.

      \includegraphics[width=6em]{../figures/copyright/by-sa}\\[2ex]
      \end{footnotesize}
        
        \begin{scriptsize}
        Typeset in Lucida Sans, \textrm{Luicda Bright}, \texttt{Lucida Console} and Lucida Math. Icons from fonts ``WebHostingHub Glyphs'' (under SIL-Open Font License) from \url{https://www.webhostinghub.com/}; ``insect icons'' (free from \url{http://www.woodcutter.es/}); ``linea-basic-10'' and ``linea-weather-10'' (free from \url{https://github.com/linea-io}), ``Mini Pics Uprooted Twig'' and ``Mini Pics Uprooted Twig'' (commercial, from Image Club Graphics, Inc.). Plant icon as .svg by Abdul Wahhab (free from \url{NounProject.com}).

        Illustrations and text quoted from copyrighted sources is excluded from this license and their use should respect the original licenses.
        \end{scriptsize}
    \end{center}
  \end{frame}


  \begin{frame}
    \frametitle{Outline}
    \tableofcontents
  \end{frame}

\section{Experiments and surveys}

\begin{frame}{Types of experiments (refresher)}
\begin{small}
    \begin{description}[type=1]
        \item[Manipulative experiments.] With manipulative experiments it
        is possible to \emph{directly} demonstrate cause-effect relationships.

        \item[Observational experiments.] When manipulative experiments
        are not possible, direct demonstration of cause-effect
        relationships can be very difficult.

        \item[Impossible experiments.] In many situations
        manipulative experiments are physically or ethically
        impossible to carry out.

        \item[Simulations.] When we do experiments with models
        instead of the real system we call them simulations.

        \item[\DExamples]

    \end{description}
\end{small}
\end{frame}

\begin{frame}{Experiments and causality}
\begin{itemize}
\item Demonstration of causes. In manipulative experiments we (attempt to) keep
everything equal among treatments except for the experimental
variable whose effect we want to test.

\item In observational studies we cannot be sure of which things
are different among the different situations (``treatments''), so we
can make guesses about causes, but can seldom directly demonstrate them.

\item In observational studies we can use path analysis to indirectly strengthen the evidence for causation (we add an additional dimension, usually time).

\item \DExamples

\end{itemize}
\end{frame}

\begin{frame}{Cause vs.\ correlation}

Tight correlation does not demonstrate causality, two responses may
be dependent on a third variable and be correlated without one being
the cause of the other. \vspace{5mm} \pause

\emph{Example:} correlation between number of swimmers drowning in
lakes, and sales of ice cream. Both peak in the summer, but there is
no cause-effect relationship between them.

\emph{Possible way out:} path analysis. E.g.\ have dead swimmers eaten
ice cream just before drowning? (more strictly, \ldots more frequently than survivors).

\end{frame}

\begin{frame}{Cause and effect in Biology \Discussion 10 + 5 min}

\begin{centering}
$\to$ \emph{fitness} $\to$ genome $\to$  epigenome $\to$ phenome $\to$ \emph{fitness} $\to \cdots$
\vspace{3ex}
\end{centering}

\begin{description}
%  \item[time sequence] .
  \item[\DiscussionI] Which transitions, shown as $\to$, depend directly on the environment?
  \item[\DiscussionI] ``Evolution determines physiology'' or ``Physiology determines fitness''?
  \item[\Attention] The course of evolution is not the result of natural selection alone!
\end{description}
\pause
\Attention \emph{Multiple viewpoints are possible, and not necessarily contradictory!}
\end{frame}


  \section*{References}
  \begin{frame}[t,allowframebreaks]
    \frametitle{References}
    \printbibliography
  \end{frame}

\end{document}
